\clearpage
\section*{Appendix R\quad Computational Receipts}\label{app:receipts}
\addcontentsline{toc}{section}{Appendix R: Computational Receipts}
\small\noindent Each computation marked \rcptmarker\ in the text is verified by a runnable receipt in the \texttt{receipts/} directory that ships with this work; status \textsf{OK} = verified (traced, run, output matches the claim). This appendix is generated from the receipt index, not maintained by hand.\par\medskip
\begingroup\renewcommand{\arraystretch}{1.25}
\begin{description}
\item[\rcptlabel{P08_matter_functional}{P8R1}\textbf{P8R1}\quad\texttt{P08\_matter\_functional}]\hfill\textsf{[OK]}\\
\textit{sec:gauge/thm:kernel/prop:bend} \ (matter functional 8pi T\textasciicircum{}t\_t=(r f'+f-1)/r\textasciicircum{}2+Lambda (Einstein tensor derived); lock p\_r=-rho; vacuum kernel = SdS (general soln); density rho=m'/(4 pi r\textasciicircum{}2) is the bend). 
\emph{Computes:} full Einstein tensor of ds\textasciicircum{}2=-f dt\textasciicircum{}2+dr\textasciicircum{}2/f+r\textasciicircum{}2dOmega\textasciicircum{}2 derived; eq:Ttt/eq:Ttheta matched; dsolve gives SdS; density=bend; control straight vs bent
 \emph{Bound:} covers 3 P8 claims
\ \emph{Run:} \texttt{python3 receipts/P08\_slicing\_operator/P08\_matter\_functional.py}
\item[\rcptlabel{P08_lapse_split}{P8R2}\textbf{P8R2}\quad\texttt{P08\_lapse\_split}]\hfill\textsf{[OK]}\\
\textit{prop:lapse} \ (**!! MARKED c54.161: the test that G\textasciicircum{}t\_t is independent of A is partly vacuous** -- `A not in Gtt.free\_symbols` is always true because A is an applied Function and free\_symbols holds only Symbols. Only the sibling has() test has teeth. two-function metric ds\textasciicircum{}2=-A dt\textasciicircum{}2+dr\textasciicircum{}2/f+r\textasciicircum{}2dOmega\textasciicircum{}2: density depends on f alone (indep of A); radial EoS 8pi(p\_r+rho)=(f/r)d/dr ln(A/f); lock A=f gives p\_r=-rho). 
\emph{Computes:} two-function Einstein tensor derived; G\textasciicircum{}t\_t indep of A; EoS matched; lock kills the log; control A=f*g -> (f/r)g'/g
 \emph{Bound:} ---
\ \emph{Run:} \texttt{python3 receipts/P08\_slicing\_operator/P08\_lapse\_split.py}
\item[\rcptlabel{P08_E1_cosmology}{P8R3}\textbf{P8R3}\quad\texttt{P08\_E1\_cosmology}]\hfill\textsf{[OK]}\\
\textit{prop:cosmo} \ (E=1 marginally-bound SdS geodesic = flat-LCDM cosmology: (dr/dtau)\textasciicircum{}2=2M/r+r\textasciicircum{}2/a\textasciicircum{}2 solved by (2M a\textasciicircum{}2)\textasciicircum{}\{1/3\}sinh\textasciicircum{}\{2/3\}(3tau/2a); Friedmann H\textasciicircum{}2=Lambda/3+8pi rho/3; OS at Lambda->0; HEB handover at r\_star). 
\emph{Computes:} geodesic derived (dr/dtau)\textasciicircum{}2=E\textasciicircum{}2-f; scale factor solves eq:E1; coth\textasciicircum{}2=1+csch\textasciicircum{}2 split; Lambda->0 -> tau\textasciicircum{}\{2/3\}; min of 2M/r+r\textasciicircum{}2/a\textasciicircum{}2 at r\_star=(M a\textasciicircum{}2)\textasciicircum{}\{1/3\} (HEB); control E<1 turns around
 \emph{Bound:} ties prop:cosmo to the Hubble-Eddington radius
\ \emph{Run:} \texttt{python3 receipts/P08\_slicing\_operator/P08\_E1\_cosmology.py}
\item[\rcptlabel{P08_nariai_signed_roots}{P8R4}\textbf{P8R4}\quad\texttt{P08\_nariai\_signed\_roots}]\hfill\textsf{[OK]}\\
\textit{sec:open} \ (the forced scale 2/sqrt3: FOUR distinct quantities share the magnitude (Nariai offset \textbar{}r\_0\textbar{}, ellipse focal distance, throat-image radius, lone/backward-radial root) and only the last is a signed root; on the M>0 member the pair is degenerate at +1/sqrt3 and the lone root is -2/sqrt3). 
\emph{Computes:} Nariai masses +/-2/(3sqrt3) via discriminant; signed roots at both members; R-conjugacy of the two sets; degeneracy of the pair; r\_0 always itself a root, and which offset gives 2M>0; 'larger' true in magnitude and false in value; the four same-magnitude quantities separated
 \emph{Bound:} the vantage-dependence of the signs is P2 sec:ring's, imported not re-established
\ \emph{Run:} \texttt{python3 receipts/P08\_slicing\_operator/P08\_nariai\_signed\_roots.py}
\item[\rcptlabel{P08_trichotomy}{P8R5}\textbf{P8R5}\quad\texttt{P08\_trichotomy}]\hfill\textsf{[OK]}\\
\textit{prop:trichotomy} \ (**!! MARKED c54.161: "eq:Ek verified at E=1 and E=sqrt2" is the constant arithmetic (1-1)==0 and (2-1)==1, and the flat-leaf CONTROL is simplify(6*0/a**2)==0** -- a constant zero compared with zero. Neither conjunct depends on the derived geodesic, so the energy-curvature relation was never tested at either energy and the control cannot distinguish the flat leaf from anything. three constant-curvature leaves (closed S\textasciicircum{}3 k=+1, flat horosphere k=0, open H\textasciicircum{}3 k=-1) = one SdS congruence at three energies via -k=E\textasciicircum{}2-1; Friedmann curvature term = \textasciicircum{}3R/6=6k/a\textasciicircum{}2). 
\emph{Computes:} open+closed induced metrics DERIVED from 5D embeddings (sinh/cosh de Sitter slicings); \textasciicircum{}3R computed (+6/-6) => 6k/a\textasciicircum{}2; -k=E\textasciicircum{}2-1 off the geodesic
 \emph{Bound:} closes P8
\ \emph{Run:} \texttt{python3 receipts/P08\_slicing\_operator/P08\_trichotomy.py}
\item[\rcptlabel{P08_synchronous_horosphere}{P8R6}\textbf{P8R6}\quad\texttt{P08\_synchronous\_horosphere}]\hfill\textsf{[OK]}\\
\textit{sec:synchronous} \ (flat synchronous slicing = horosphere family between two null directions: embedding on hyperboloid; induced -dtau\textasciicircum{}2+e\textasciicircum{}\{2tau/a\}dx\textasciicircum{}2 (flat k=0); eta(X,B)=(a\textasciicircum{}2/2)e\textasciicircum{}\{tau/a\} labels slices; tau->-inf singularity). 
\emph{Computes:} flat-slicing embedding + induced metric DERIVED; A,B null (eta(A,B)=a\textasciicircum{}2/2); eta(X,B) prop e\textasciicircum{}\{tau/a\} indep of transverse; tau->-inf limit 0
 \emph{Bound:} the 'direct computation' claim (not a numbered prop)
\ \emph{Run:} \texttt{python3 receipts/P08\_slicing\_operator/P08\_synchronous\_horosphere.py}
\item[\rcptlabel{P08_the_second_ruling_reading_does_not_extend_to_the_collapse_face}{P8R7}\textbf{P8R7}\quad\texttt{P08\_the\_second\_ruling\_reading\_does\_not\_extend\_to\_the\_collapse\_face}]\hfill\textsf{[ALL PASS]}\\
\textit{sec:synchronous} \ (**THE SECOND-RULING READING DOES NOT EXTEND TO THE COLLAPSE FACE, AND THE OBSTRUCTION IS THE ORTHOGONALITY THE CONSTRUCTION ITSELF REQUIRES.** P08 builds the synchronous slices as the horospheres NORMAL to the past generator B; a collapse horizon's limiting direction is generically NON-orthogonal to any spacelike slice (P01/F1) and meets its horizon tangentially (P07). **The requirement and its failure are the same property.** The two faces are separated by the horizon cubic's discriminant, not a modelling choice: massless gives 4 alpha\textasciicircum{}6 with three distinct roots and a transverse crossing, the forced member vanishes on the merged root at the front seam.). 
\emph{Computes:} r4239
 \emph{Bound:} 61
\ \emph{Run:} \texttt{python3 receipts/P08\_slicing\_operator/P08\_the\_second\_ruling\_reading\_does\_not\_extend\_to\_the\_collapse\_face.py}
\item[\rcptlabel{P08_the_closure_is_the_constitutive_relation}{P8R8}\textbf{P8R8}\quad\texttt{P08\_the\_closure\_is\_the\_constitutive\_relation}]\hfill\textsf{[OK]}\\
\textit{sec:open} \ (the spherical class CLOSES: the contracted Bianchi identity is entailed by the cut's own geometry and fixes the angular component from the radial pair, and a constitutive relation closes the system to the TOV pair in the operator's own variables -- so what the construction does not supply is the equation of state, which general relativity does not supply either). 
\emph{Computes:} symbolic on the general static spherical cut; load-bearing NEGATIVES: a non-zero Bianchi residual breaks the count, a lapse equation not reducing to Schwarzschild breaks the closure
 \emph{Bound:} the equation of state itself is NOT established and is the input, here and in general relativity alike
\ \emph{Run:} \texttt{python3 receipts/P08\_slicing\_operator/P08\_the\_closure\_is\_the\_constitutive\_relation.py}
\item[\rcptlabel{P08_the_content_law_and_the_head_are_one_wall_seen_twice}{P8R9}\textbf{P8R9}\quad\texttt{P08\_the\_content\_law\_and\_the\_head\_are\_one\_wall\_seen\_twice}]\hfill\textsf{[OK]}\\
\textit{the rooms map (`r6443`); `PO-30`, `PO-31`, and `PO-41`'s strike} \ (**rooms (2) and (3) are one wall seen twice, and the map drew them apart.** `PO-30` wants a **law** (why a cut bends); the head owes a **fact** (baryons on its collapse leg, and how many). Every lap after the first inherits its cut's shape and ***the head inherits nothing***, so the law that fixes a cut absent a progenitor is exactly what would say what matter the head carries. \ensuremath{\Rightarrow} **`PO-30` discharged discharges both; the converse fails**). 
\emph{Computes:} what each row owes by kind, the head's non-inheritance (`PO-40`'s result), and the entailment tested in both directions
 \emph{Bound:}  **a structural reading, not a computation** --- it fails if `PO-30`'s law fixes only the cut-content RELATION and not which cut, and that is named as where to attack it
\ \emph{Run:} \texttt{python3 receipts/P08\_slicing\_operator/P08\_the\_content\_law\_and\_the\_head\_are\_one\_wall\_seen\_twice.py}
\item[\rcptlabel{P08_the_one_distinguished_cut_carries_no_matter}{P8R10}\textbf{P8R10}\quad\texttt{P08\_the\_one\_distinguished\_cut\_carries\_no\_matter}]\hfill\textsf{[OK]}\\
\textit{`PO-30`; `p0`'s offset-mass relation; `PO-41`'s head} \ (**the offset-mass cubic \$2M=\textbackslash\{\}alpha(x-x\textasciicircum{}\{3\})\$ has a unique positive stationary point and it is EXACTLY the Nariai member** --- \$r\_0=\textbackslash\{\}alpha/\textbackslash\{\}sqrt3=r\_N\$, \$M=\textbackslash\{\}alpha\textbackslash\{\}sqrt3/9=M\_N\$ --- ***a third route to Nariai, by extremising, where `P05` reaches it as \$\textbackslash\{\}sigma\$'s fixed point and `P07` by a trichotomy; unremarked in the corpus at any of the four sites the cubic appears***. **But every member of that family has \$m(r)\$ CONSTANT, so \$\textbackslash\{\}rho=m'/4\textbackslash\{\}pi r\textasciicircum{}\{2\}=0\$ identically** --- the one intrinsically distinguished cut carries **no distributed matter**, where the head owes baryons. \ensuremath{\Rightarrow} **the law must fix a FUNCTION \$m(r)\$, not a number**). 
\emph{Computes:} the stationary point solved and checked against the corpus's own \$r\_N\$ and \$M\_N\$, the second derivative confirming a maximum, and \$\textbackslash\{\}rho=0\$ on a constant offset
 \emph{Bound:} **not claimed that no such law exists** --- only that it cannot be *take the distinguished member*, that member being empty; and Nariai remains forced on the **collapse** by `P07`, which is a statement about a limit and not about the profile that collapses
\ \emph{Run:} \texttt{python3 receipts/P08\_slicing\_operator/P08\_the\_one\_distinguished\_cut\_carries\_no\_matter.py}
\item[\rcptlabel{P08_the_selection_principle_cannot_supply_matter}{P8R11}\textbf{P8R11}\quad\texttt{P08\_the\_selection\_principle\_cannot\_supply\_matter}]\hfill\textsf{[OK]}\\
\textit{`PO-30`; `P12`'s least-arbitrariness, `p0`'s constant-density clause, `P07`'s wall} \ (**the corpus's only selection principle cannot be the law this row wants, and the reason is structural.** `P12`: it prefers the structure requiring **no choice of how to break a symmetry** --- among profiles, exactly one qualifies, \$\textbackslash\{\}rho\$ constant. **And `p0` says a constant density "enters the profile's \$\textbackslash\{\}Lambda r\textasciicircum{}\{2\}/3\$ term, NOT as a \$2m/r\$ bend"** --- so the selected profile *is the cosmological term, not matter*; `P07` reaches the same boundary from the other side, homogeneous collapse lying past the wall where generation-by-symmetry hands off to Einstein evolution. \ensuremath{\Rightarrow} ***matter is a BREAKING of maximal symmetry and the principle selects what requires no breaking: it cannot supply it***). 
\emph{Computes:} the profile enumeration against `P12`'s own criterion, the term a constant density enters per `p0`, and the selector/thing incompatibility
 \emph{Bound:}  **it does not close the row** --- it removes the one existing candidate, which narrows what an answer could look like; and it does not claim no such law exists
\ \emph{Run:} \texttt{python3 receipts/P08\_slicing\_operator/P08\_the\_selection\_principle\_cannot\_supply\_matter.py}
\item[\rcptlabel{P08_the_construction_already_generates_content_and_how}{P8R12}\textbf{P8R12}\quad\texttt{P08\_the\_construction\_already\_generates\_content\_and\_how}]\hfill\textsf{[OK]}\\
\textit{`PO-30`; `P14`'s `prop:forced`, `p0`'s discrete residue} \ (**the criterion chooses AMONG BREAKINGS and takes the one that is itself symmetric** --- `P14`: a one-plane construction "must select **which** hinge: a free modulus", and the \$\textbackslash\{\}mathbb\{Z\}\_3\$-symmetric three-plane is "the unique configuration carrying no such modulus". \ensuremath{\Rightarrow} ***And by that route `P14` already delivers content --- generation count, chirality, family symmetry*** --- so content **is** generated here. And the wall, the generations and the chirality are **one residue read at three places**: \$\textbackslash\{\}mathrm\{Aut\}(A\_2)=S\_3\textbackslash\{\}times\textbackslash\{\}mathbb\{Z\}\_2\$ off the waist, \$S\_3\$ on the points **on** the circle and \$\textbackslash\{\}mathbb\{Z\}\_2\$ on the lines **tangent** to it). 
\emph{Computes:} the two available configurations against `P14`'s own criterion, the content that route delivers, the residue's factorisation, and the discrete/continuous asymmetry
 \emph{Bound:}  **closes nothing.** What it establishes is that the route is real, has produced content once, and what the open half is: **whether the CONTINUOUS content has a moduli-free configuration as the discrete content does** --- not claimed that it has
\ \emph{Run:} \texttt{python3 receipts/P08\_slicing\_operator/P08\_the\_construction\_already\_generates\_content\_and\_how.py}
\item[\rcptlabel{K6_fullness}{P8R13}\textbf{P8R13}\quad\texttt{K6\_fullness}]\hfill\textsf{[OK]}\\
\textit{eq:gauge} \ (**!! MARKED c54.161: the stratum-by-stratum match column compares each row's dimension WITH ITSELF** -- both fields are the same expression written twice, so the column can never print NO and the fullness claim went unchecked. The added assertions pin the dimensions against P12's independently. K6 --- IS Psi FULL? P5's own stated gap, posed functorially by K2, and P12 already has the data. CONFIRMATION 2 -- read P12's stratification statement WHOLE, which a first pass did not: 'the symmetric-pair isotropy dimensions of so(5,1) are \{6,7,10\}, so Type O (SO(4,1), ten) and Nariai (SO(2,1)xSO(3), six) sit at the only admissible dimensions, while EVERY GENERIC STRATUM (FOUR, three, two, zero) is non-symmetric by di...). 
\emph{Computes:} runs clean; registered from its own docstring and a live run
 \emph{Bound:} description is the receipt's own wording, condensed; not independently re-derived here
\ \emph{Run:} \texttt{python3 receipts/P08\_slicing\_operator/K6\_fullness.py}
\item[\rcptlabel{K9_isotropy_obstruction}{P8R14}\textbf{P8R14}\quad\texttt{K9\_isotropy\_obstruction}]\hfill\textsf{[OK]}\\
\textit{eq:gauge} \ (K9 --- IS THE ISOTROPY=ISOMETRY EQUALITY GENERAL, OR STRATUM-BY-STRATUM? I hedged this twice (K6 in P8, K7 in P9) and proposed the hedge could be removed by a general argument: 'a cut's isotropy is by construction the subgroup preserving it, and the geometry read off it inherits exactly those.' TEST THAT ARGUMENT HONESTLY, BOTH DIRECTIONS.). 
\emph{Computes:} runs clean; registered from its own docstring and a live run
 \emph{Bound:} description is the receipt's own wording, condensed; not independently re-derived here
\ \emph{Run:} \texttt{python3 receipts/P08\_slicing\_operator/K9\_isotropy\_obstruction.py}
\item[\rcptlabel{V34_orthogonal_and_no_cost}{P8R15}\textbf{P8R15}\quad\texttt{V34\_orthogonal\_and\_no\_cost}]\hfill\textsf{[OK]}\\
\textit{eq:gauge} \ (V3 --- IS THE LEFTWARD READING INCOMPATIBLE WITH AN ACTION PRINCIPLE, OR ORTHOGONAL TO ONE? V4 --- WHAT DOES TAKING THE DIRAC ALGEBRA WITHOUT A LEGENDRE TRANSFORM COST? CONFIRMATION 2 first, and it moves both answers.). 
\emph{Computes:} runs clean; registered from its own docstring and a live run
 \emph{Bound:} description is the receipt's own wording, condensed; not independently re-derived here
\ \emph{Run:} \texttt{python3 receipts/P08\_slicing\_operator/V34\_orthogonal\_and\_no\_cost.py}
\item[\rcptlabel{P08_the_branch_point_monodromy_is_Z3}{P8R16}\textbf{P8R16}\quad\texttt{P08\_the\_branch\_point\_monodromy\_is\_Z3}]\hfill\textsf{[OK]}\\
\textit{prop:cosmo} \ (**THE COSMOGENESIS BRANCH POINT CARRIES A Z\_3 MONODROMY, AND ONLY TWO OF ITS THREE SHEETS CARRY REAL GEOMETRY.** In r = (2 M alpha\textasciicircum{}2)\textasciicircum{}(1/3) sinh\textasciicircum{}(2/3)(u), u = 3 tau / 2 alpha, the sine near u = 0 is u times something analytic and non-vanishing, so **the whole branch structure is the factor u\textasciicircum{}(2/3)**: one circuit multiplies r by e\textasciicircum{}\{4 pi i/3\}, a primitive cube root of unity. **One circuit does not close it; three do.** ** OF THE THREE DETERMINATIONS EXACTLY TWO CARRY REAL r, and they are the two the construction already uses: ** the real u axis gives r > 0 (the expanding leg) and Im u = -pi/2 -- the offset Im tautilde = -pi alpha/3 the companion papers continue along -- gives r = -(2M alpha\textasciicircum{}2)\textasciicircum{}(1/3) cosh\textasciicircum{}(2/3), the conjugate branch. **The third sheet carries no real geometry to continue onto.** ** AND THE ORDER IS THE MASS TERM'S, NOT THE HORIZON CUBIC'S: ** near r -> 0 the 2M/r in (dr/dtau)\textasciicircum{}2 dominates and gives r\textasciicircum{}(3/2) \textasciitilde{} tau, whereas the cubic's degree comes from the r\textasciicircum{}2/alpha\textasciicircum{}2 term. **Two threes of the same order and different origin, recorded as distinct.**). 
\emph{Computes:} the monodromy multiplier asserted a primitive cube root of unity with order asserted 3; the real-r sheets tabulated line by line; the horizon cubic's degree asserted 3 with its leading term identified; and R asserted central in the Weyl S\_3, an involution, and outside it
 \emph{Bound:} ** THIS IS A THIRD THREE AND IS A LIABILITY BEFORE IT IS AN ASSET. ** The programme has already withdrawn one conflation of two threes (generations-are-the-hinges, registered in `check\_withdrawn`), so this one enters labelled as DISTINCT from the A\_2 / horizon-root / hinge three. What would make them one is an argument that the 2/3 exponent and the cubic's degree are forced by a single structure; **none is offered**
\ \emph{Run:} \texttt{python3 receipts/P08\_slicing\_operator/P08\_the\_branch\_point\_monodromy\_is\_Z3.py}
\item[\rcptlabel{I12_the_marginal_member_is_the_only_non_degenerate_one_that_integrates}{P8R17}\textbf{P8R17}\quad\texttt{I12\_the\_marginal\_member\_is\_the\_only\_non\_degenerate\_one\_that\_integrates}]\hfill\textsf{[OK]}\\
\textit{sec:trichotomy} \ (I12 --- THE MARGINAL MEMBER IS THE ONLY NON-DEGENERATE ONE THAT INTEGRATES. Integrable-systems field bake, probe I12. Establishes that the E=1 quadrature is elementary and inverts to sinh\textasciicircum{}(2/3); that at k != 0 the substitution r = u\textasciicircum{}2 gives a SEXTIC under the root and the integral returns unevaluated; that the only other elementary cases are the degenerate corners Lambda=0 (the cycloid) and M=0 (pure de Sitter); and that the cycloid satisfies the Lambda=0 law identically and NOT P08's own closed-member law.). 
\emph{Computes:} runs clean; ten verdicts
 \emph{Bound:} COMPUTES: the dropped term's weight r\textasciicircum{}3/2M alpha\textasciicircum{}2 is evaluated at the illustrative gauge 2M = 0.1 alpha, giving 1e-5 at r/alpha=0.01 and 0.27 at 0.3 --- an illustration of the ratio's growth, not a physical parameter of any member
\ \emph{Run:} \texttt{python3 receipts/P08\_slicing\_operator/I12\_the\_marginal\_member\_is\_the\_only\_non\_degenerate\_one\_that\_integrates.py}
\item[\label{rcpt:W1_what_remains_between_the_wall_and_a_curve_dynamics}\textbf{}\quad\texttt{W1\_what\_remains\_between\_the\_wall\_and\_a\_curve\_dynamics}]\hfill\textsf{[OK]}\\
\textit{sec:open} \ (the general inhomogeneous evolution is ordinary leaf evolution, not a new generative law). 
\emph{Computes:} asserts P8 kinematic split, the Hamiltonian-constraint form, and that the raises phrasing is comment-only (13)
 \emph{Bound:} NOT-A-PAPER-CLAIM --- discharges L-207 section 1
\ \emph{Run:} \texttt{python3 receipts/L207\_the\_bend/W1\_what\_remains\_between\_the\_wall\_and\_a\_curve\_dynamics.py}
\item[\label{rcpt:W2_the_general_inhomogeneous_leaf_exhibited}\textbf{}\quad\texttt{W2\_the\_general\_inhomogeneous\_leaf\_exhibited}]\hfill\textsf{[OK]}\\
\textit{sec:open} \ (the general inhomogeneous leaf exhibited: LTB with arbitrary m(r), the bend-density identity exact and the evolution ordinary GR). 
\emph{Computes:} computes G\_tt from the LTB metric and asserts the identity and the acceleration law symbolically (11 checks)
 \emph{Bound:} NOT-A-PAPER-CLAIM --- discharges L-207 (1)
\ \emph{Run:} \texttt{python3 receipts/L207\_the\_bend/W2\_the\_general\_inhomogeneous\_leaf\_exhibited.py}
\item[\rcptlabel{P03_acceleration_is_slice_curvature}{P3R12}\textbf{P3R12}\quad\texttt{P03\_acceleration\_is\_slice\_curvature}]\hfill\textsf{[OK]}\\
\textit{sec:curvature \ensuremath{\cdot} rem:twocritical \ensuremath{\cdot} the bend} \ (on ANY member of the energy family (E cancels: (dr/dtau)\textasciicircum{}2 = E\textasciicircum{}2 - f, and E is a constant of the motion) d\textasciicircum{}2r/dtau\textasciicircum{}2 = -f'/2 = r*K\_G identically, so THE SIGN OF THE SLICE'S GAUSSIAN CURVATURE IS THE SIGN OF THE COMOVING ACCELERATION -- and the flat locus is the deceleration-to-acceleration turnover, which on Nariai is the seam and in standard terms is rho\_m = 2 rho\_Lambda (csch\textasciicircum{}2 w = 2 <=> sinh w = 1/sqrt2, the seam's phase). A recent cosmological epoch: 7.06 Gyr at H0=73, 7.65 at 67.4, preceding matter-Lambda equality [borrowed from P3 \ensuremath{\cdot} P7 \ensuremath{\cdot} P8]). 
\emph{Computes:} the identity asserted; the K\_G zero solved and matched to (M alpha\textasciicircum{}2)\textasciicircum{}(1/3); csch\textasciicircum{}2 w = 2 <=> sinh w = 1/sqrt2 proved by algebra after simplify() failed on the closed forms; epochs at both H0
 \emph{Bound:} bound: identities on the Nariai E=1 worldline; no causal claim between seam and turnover -- they are one locus; dating inherits H0
\ \emph{Run:} \texttt{python3 receipts/P03\_SdS\_slicing/P03\_acceleration\_is\_slice\_curvature.py}
\item[\rcptlabel{order3_bridge}{P7R6}\textbf{P7R6}\quad\texttt{order3\_bridge}]\hfill\textsf{[OK]}\\
\textit{groupoid<->bead order-three} \ (the groupoid Z/3 (horizon cubic r\textasciicircum{}3-r+2M, colinear real) and bead Z/3 (comoving r\textasciicircum{}3+2M, equilateral) are Z/3's of DIFFERENT cubics and no AFFINE change of variable identifies their root sets (obstruction at fixed cubic; the two are the E=1 and E=0 ends of one turning-point family, stage 5) [borrowed from P7]). 
\emph{Computes:} bead law obeys (dr/dtau\textasciitilde{})\textasciicircum{}2=1-f; horizon roots colinear-real vs comoving equilateral; not affinely equivalent; stage 5 exhibits the E-family joining them, with a discriminating control
 \emph{Bound:} the affine obstruction is exact; identification of ROOT SETS remains unclaimed
\ \emph{Run:} \texttt{python3 receipts/P07\_CR\_framework/order3\_bridge.py}
\item[\rcptlabel{K23_functor_and_quotient}{P5R8}\textbf{P5R8}\quad\texttt{K23\_functor\_and\_quotient}]\hfill\textsf{[OK]}\\
\textit{rem:universal} \ (K2 + K3 --- IS Psi A FUNCTOR, AND IS 'THE SINGLE INVARIANT IS THE GEOMETRY' A UNIVERSAL PROPERTY? CONFIRMATION 4 --- the objects: SOURCE: K1's action groupoid SO(5,1) \textbar{}x C ==> C . Objects = cuts. Arrows c -> g.c . Psi (P8): 'carries a cut of the substrate to a geometry'. TARGET: geometries, with isometries between them. [borrowed from P5 (+P8)]). 
\emph{Computes:} runs clean; registered from its own docstring and a live run
 \emph{Bound:} description is the receipt's own wording, condensed; not independently re-derived here
\ \emph{Run:} \texttt{python3 receipts/P05\_groupoid/K23\_functor\_and\_quotient.py}
\item[\rcptlabel{I52_the_vacuum_theorem_is_an_index_one_statement_and_the_index_is_the_mass}{P8R18}\textbf{P8R18}\quad\texttt{I52\_the\_vacuum\_theorem\_is\_an\_index\_one\_statement\_and\_the\_index\_is\_the\_mass}]\hfill\textsf{[OK]}\\
\textit{sec:kernel, sec:bend} \ (P08 states dim ker = 1 ("the single constant of integration -2M") and proves the matter functional onto ("matter is bend": any density is realised by some cut), and never assembles them: `kernel` x14, `onto` x18, `surjectiv` x3, `cokernel` x0, `Fredholm` x0. Those two numbers ARE a Fredholm index, 1 - 0 = 1, and the one dimension of kernel IS the mass [borrowed from P08]). 
\emph{Computes:} kernel dimension by INTEGRATION (solutions from independent initial data, rank of the sampled matrix) returns 1, with f''=0 as the operator-varying control returning 2; cokernel exhibited empty by CONSTRUCTING a preimage f=(1/r)int g for 40 random forcings, worst residual 3.4e-06
 \emph{Bound:} SUPPORTS P08 and names what it has --- an index is deformation-stable, so this is stronger than the theorem's "this ODE has a one-parameter solution set"
\ \emph{Run:} \texttt{python3 receipts/P08\_slicing\_operator/I52\_the\_vacuum\_theorem\_is\_an\_index\_one\_statement\_and\_the\_index\_is\_the\_mass.py}
\end{description}\endgroup
